Die Gewichtskraft greift im Schwerpunkt S an und zeigt senkrecht nach unten. Die Kraft der Achse greift in D an und hat bezüglich D kein Drehmoment. Das Brett bleibt also nur in Ruhe, wenn auch die Gewichtskraft kein Drehmoment hat, d.\,h. wenn ihre Wirkungslinie, die Senkrechte durch S, durch D geht. Ob das Gleichgewicht stabil ist, zeigt ein Gedankenexperiment: Man dreht das Brett ein wenig und schaut, ob das Drehmoment der Gewichtskraft es zurückdreht.

\begin{abcliste}
\abc
\begin{itemize}
 \item Brett 1: S liegt senkrecht unter D, die Gewichtskraft hat keinen Hebelarm: Gleichgewicht. Dreht man das Brett ein wenig, so hebt sich S und rückt zur Seite; die Gewichtskraft dreht das Brett zurück: \result{\text{stabiles Gleichgewicht}}.
 \item Brett 2: S liegt senkrecht über D: ebenfalls Gleichgewicht. Ein wenig gedreht, senkt sich S und rückt zur Seite; die Gewichtskraft dreht das Brett weiter weg: \result{\text{labiles Gleichgewicht}}. Schon ein kleiner Stoss kippt es um.
 \item Brett 3: Die Achse geht durch S. Die Gewichtskraft greift in der Achse an und hat in jeder Lage den Hebelarm null: \result{\text{indifferentes Gleichgewicht}}. Das Brett bleibt in jeder Lage stehen.
 \item Brett 4: Die Senkrechte durch S verläuft \SI{10}{\cm} links an D vorbei. Die Gewichtskraft hat bezüglich D ein Drehmoment: \result{\text{kein Gleichgewicht}}.
\end{itemize}

\abc
Die Gewichtskraft zieht S links von D nach unten: Brett 4 beginnt sich im Gegenuhrzeigersinn zu drehen. Ohne Reibung schwingt es wie ein Pendel um die Lage, in der S senkrecht unter D liegt, und kommt bis auf die Gegenseite (S \SI{10}{\cm} rechts von D, auf gleicher Höhe). Mit etwas Reibung wird die Schwingung kleiner, und das Brett kommt in der Lage zur Ruhe, in der \result{\text{S senkrecht unter D}} liegt, um \ang{90} gedreht gegenüber der Zeichnung: ein stabiles Gleichgewicht wie bei Brett 1.
\end{abcliste}
